\documentclass[tikz,border=1.35mm]{standalone}

\definecolor{fvmadaptedge}{HTML}{8A7E7E}
\usetikzlibrary{arrows.meta}

\definecolor{splitred}{HTML}{AF2525}

% Wireframe exploded view of prism split code 2. The three separated children
% are the Prism(4) columns generated around the split triangular end faces.
\providecommand{\fvmadaptYawAngle}{8}
\providecommand{\fvmadaptExplodeAmount}{1}
\providecommand{\fvmadaptExplodeScale}{1.35}
\pgfmathsetmacro{\yawcos}{cos(\fvmadaptYawAngle)}
\pgfmathsetmacro{\yawsin}{sin(\fvmadaptYawAngle)}
\pgfmathsetmacro{\xbasisx}{20.3*\yawcos + 15.3*\yawsin}
\pgfmathsetmacro{\xbasisy}{-8.5*\yawcos + 11.3*\yawsin}
\pgfmathsetmacro{\ybasisx}{-20.3*\yawsin + 15.3*\yawcos}
\pgfmathsetmacro{\ybasisy}{8.5*\yawsin + 11.3*\yawcos}
\pgfmathsetmacro{\fvmadaptEtaProjectionScale}{0.75}
\pgfmathsetmacro{\ybasisxscaled}{\fvmadaptEtaProjectionScale*\ybasisx}
\pgfmathsetmacro{\ybasisyscaled}{\fvmadaptEtaProjectionScale*\ybasisy}
\pgfmathsetmacro{\explodeXi}{0.56*\fvmadaptExplodeScale*\fvmadaptExplodeAmount}
\pgfmathsetmacro{\explodeEta}{0.43*\fvmadaptExplodeScale*\fvmadaptExplodeAmount}

\begin{document}
\begin{tikzpicture}[
  x={(\xbasisx mm,\xbasisy mm)},
  y={(\ybasisxscaled mm,\ybasisyscaled mm)},
  z={(0mm,20mm)},
  line join=round,
  line cap=round,
  outer/.style={draw=fvmadaptedge,line width=0.54mm},
  split/.style={draw=splitred,line width=0.44mm},
  axis/.style={draw=fvmadaptedge!60,line width=0.32mm,-{Latex[length=2.5mm,width=1.8mm]}}
]
  % Fixed canvas for static and animated renders.
  \path[use as bounding box]
    (canvas cs:x=-57mm,y=-58mm) rectangle (canvas cs:x=72mm,y=68mm);

  \newcommand{\drawquadcolumn}{%
    % Bottom and top end-face fragments: original boundary pieces stay #8a7e7e,
    % while the face-center split edges are red.
    \draw[outer] (B0) -- (B1);
    \draw[split] (B1) -- (BC) -- (B3);
    \draw[outer] (B3) -- (B0);
    \draw[outer] (T0) -- (T1);
    \draw[split] (T1) -- (TC) -- (T3);
    \draw[outer] (T3) -- (T0);

    % Original vertical edge plus the new side and axial split edges.
    \draw[outer] (B0) -- (T0);
    \draw[split] (B1) -- (T1);
    \draw[split] (BC) -- (TC);
    \draw[split] (B3) -- (T3);

    \foreach \p in {B0,T0} {
      \fill[fvmadaptedge] (\p) circle[radius=0.58mm];
    }
    \foreach \p in {B1,BC,B3,T1,TC,T3} {
      \fill[splitred] (\p) circle[radius=0.52mm];
    }
  }

  % Compact orientation axes near the fixed A-column child.
  \coordinate (Oaxis) at (-1.28,-0.55,-0.70);
  \draw[axis] (Oaxis) -- (-0.32,-0.55,-0.70) node[below right,font=\normalsize,text=fvmadaptedge] {$\xi$};
  \draw[axis] (Oaxis) -- (-1.28,0.39,-0.70) node[above left,font=\normalsize,text=fvmadaptedge] {$\eta$};
  \draw[axis] (Oaxis) -- (-1.28,-0.55,0.26) node[above,font=\normalsize,text=fvmadaptedge] {$\zeta$};

  % Child around vertex A: keep the low xi/eta column fixed.
  \begin{scope}[shift={(0,0,0)}]
    \coordinate (B0) at (0,0,0);
    \coordinate (B1) at (0.725,0,0);
    \coordinate (BC) at (0.483,0.333,0);
    \coordinate (B3) at (0,0.5,0);
    \coordinate (T0) at (0,0,1);
    \coordinate (T1) at (0.725,0,1);
    \coordinate (TC) at (0.483,0.333,1);
    \coordinate (T3) at (0,0.5,1);
    \drawquadcolumn
  \end{scope}

  % Child around vertex B: shift mostly along +xi.
  \begin{scope}[shift={(\explodeXi,0,0)}]
    \coordinate (B0) at (1.45,0,0);
    \coordinate (B1) at (0.725,0.5,0);
    \coordinate (BC) at (0.483,0.333,0);
    \coordinate (B3) at (0.725,0,0);
    \coordinate (T0) at (1.45,0,1);
    \coordinate (T1) at (0.725,0.5,1);
    \coordinate (TC) at (0.483,0.333,1);
    \coordinate (T3) at (0.725,0,1);
    \drawquadcolumn
  \end{scope}

  % Child around vertex C: shift mostly along +eta.
  \begin{scope}[shift={(0,\explodeEta,0)}]
    \coordinate (B0) at (0,1,0);
    \coordinate (B1) at (0,0.5,0);
    \coordinate (BC) at (0.483,0.333,0);
    \coordinate (B3) at (0.725,0.5,0);
    \coordinate (T0) at (0,1,1);
    \coordinate (T1) at (0,0.5,1);
    \coordinate (TC) at (0.483,0.333,1);
    \coordinate (T3) at (0.725,0.5,1);
    \drawquadcolumn
  \end{scope}
\end{tikzpicture}
\end{document}
